\section{Case III: a condensed-phase WCA dimer}
\label{sec:case_wca}

The third system is many-body and is where ABF is most starved: a dimer that must
push a dense solvent aside to change its bond length. It tests whether the
two-dimensional findings survive in a setting with thousands of degrees of freedom
and a nonlinear collective variable.

\subsection{Model and reference}
A dimer of two tagged particles sits in a periodic two-dimensional bath of
$n_{\mathrm{dim}}^2-2$ Weeks--Chandler--Andersen (WCA) solvent particles
($n_{\mathrm{dim}}=10$, lattice spacing $a$, box $L=n_{\mathrm{dim}}\,a$). All
particles interact through the purely repulsive WCA
potential~\citep{WeeksChandlerAndersen1971} (Lennard--Jones cut and shifted at its
minimum $r_0=2^{1/6}\sigma$). The two dimer particles additionally feel a symmetric
double-well bond $V_{\mathrm{dim}}(r)=h\,(1-((r-r_0-w)/w)^2)^2$ with barrier
height $h=2$ and width $w=2$. The reaction coordinate is the scaled bond length
\begin{equation}
  z(q) = \frac{|q_1-q_2|-r_0}{2w},
  \label{eq:wca_cv}
\end{equation}
a nonlinear collective variable of the full configuration $q\in\R^{2n}$
(\cref{rem:manybody}): $z\approx0$ is the \emph{compact} state, $z\approx1$ the
\emph{stretched} state, and $z\in[0.25,0.75]$ the \emph{transition} region. The
free-energy barrier near $z=0.5$ is partly entropic / crowding-driven, since the
dimer must displace solvent to cross it. Dynamics are overdamped Langevin at
$\beta=1$, $\mathrm{d}t=0.002$, with the minimum-image convention; the default
system is $a=1.5$, $\sigma=\epsilon=1$. The local mean force along $z$ is the
constrained projection of the physical forces plus the geometric $2w/(\beta r)$
term (\cref{rem:manybody}). The reference $F(z)$ is computed once per system $a$ by
constrained-MD thermodynamic integration~\citep{Carter1989,DenOtterBriels1998}
and used for evaluation only; errors are taken on the interior window
$z\in[-0.1,1.1]$ with the additive constant aligned.

\subsection{Methods and design}
As before the practical comparison is \textbf{ABF only} versus \textbf{ABF+FR
estimated}~\eqref{eq:q_est}, with \textbf{FR uniform} and \textbf{FR oracle}
diagnostics; here the oracle target is built from the TI reference and is
explicitly \emph{not} deployable (\cref{rem:oracle}). A no-leakage assertion
ensures only the oracle method ever receives the reference, and a clone copies the
\emph{entire} WCA configuration --- not just $z$ --- so cloned solvent
arrangements travel with the bond length. Replica lineage is tracked to report the
ancestor effective sample size $\mathrm{ESS}=1/\sum_a w_a^2$. The tuned FR
configuration is given in \cref{tab:wca_design}; the rate is named
\texttt{fr\_rate} and the onset \texttt{fr\_start\_steps} (\cref{sec:fr_notation}).

\begin{table}[t]
\centering
\small
\caption{WCA dimer base system and tuned Fisher--Rao hyperparameters.}
\label{tab:wca_design}
\begin{tabular}{ll}
\toprule
Quantity & Value \\
\midrule
Inverse temperature $\beta$ & 1 \\
Time step $\Delta t$ & 0.002 \\
Steps per run & 250{,}000 \\
Replicas $N$ & 1024 \\
Lattice spacing $a$ & 1.5 \\
Dimer barrier $h$, width $w$ & 2, 2 \\
FR rate \texttt{fr\_rate} & 0.10 \\
FR onset \texttt{fr\_start\_steps} & 20{,}000 \\
FR stride \texttt{fr\_every} & 5 \\
Score clip $S_{\max}$ & 2.0 \\
Event cap $f_{\max}$ & 0.02 \\
Target EMA rate & 0.005 \\
\bottomrule
\end{tabular}
\end{table}


\subsection{Results}
The main comparison (\cref{tab:wca_main}; 10 seeds, $N=1024$, $250\,000$ steps)
shows the practical method \textbf{reliably beating ABF}: the tuned FR cuts the
median final $\Ltwo(F)$ from $\WCAabfLtwoF$ to $\WCAtunedLtwoF$
(\,$+\WCAtunedGainPct\%$\,), winning all $\WCAtunedWin/10$ matched seeds. As in the
entropic-bottleneck case, the uniform target ($+\WCAuniGainPct\%$) and the oracle
target ($+\WCAoracleGainPct\%$) are \emph{essentially as good as} the estimated
target --- the online estimated target already extracts nearly all the available
benefit, and there is no headroom to the oracle. We are explicit that the headline
gain is in the \emph{free energy}: the mean-force error $\Ltwo(F')$ improves only
modestly ($\approx22\%$), so we do not claim a large mean-force improvement.

\begin{table}[t]
\centering
\small
\caption{WCA main comparison (250k steps, $N=1024$, $a=1.5$). The headline gain is in the free energy; the mean-force gain is modest.}
\label{tab:wca_main}
\begin{tabular}{lrrrrr}
\toprule
Method & $\|F\!-\!F_{\rm ref}\|$ & $\|F'\!-\!F'_{\rm ref}\|$ & gain (\%) & wins & anc. ESS \\
\midrule
ABF only & 0.0835 & 0.3220 & -- & -- & -- \\
FR estimated (tuned) & 0.0361 & 0.2513 & +55.0 & 10/10 & 86 \\
FR uniform & 0.0372 & 0.2525 & +54.6 & 10/10 & 90 \\
FR oracle$^{*}$ & 0.0387 & 0.2475 & +52.7 & 10/10 & 100 \\
FR estimated (strong) & 0.0329 & 0.2590 & +58.9 & 10/10 & 45 \\
FR estimated (aggressive) & 0.0935 & 0.3883 & -13.1 & 1/10 & 22 \\
\bottomrule
\end{tabular}
\par\smallskip\footnotesize $^{*}$\,The oracle target uses the thermodynamic-integration reference $F_{\rm ref}$; it is a diagnostic control, \emph{not} a deployable method. Medians over $10$ seeds.
\end{table}


\begin{figure}[t]
  \centering
  \includegraphics[width=0.92\textwidth]{fig_wca_01_convergence.png}
  \caption{WCA dimer, convergence (median and inter-quartile range over seeds).
    ABF and FR are identical until FR activates; from then on FR is consistently
    below ABF on $\Ltwo(F)$ for the rest of the run. The cumulative replacement
    count (right) confirms the birth--death is active but bounded.}
  \label{fig:wca_convergence}
\end{figure}

\begin{figure}[t]
  \centering
  \includegraphics[width=0.92\textwidth]{fig_wca_02_profiles.png}
  \caption{Final profiles (median over seeds) versus the TI reference (dashed):
    (a) free energy $\Fhat(z)$, (b) mean force $\Fphat(z)$, (c) reaction-coordinate
    marginal. FR recovers the reference free energy more accurately than ABF,
    chiefly in and around the transition region.}
  \label{fig:wca_profiles}
\end{figure}

\subsection{Mechanism: resampling-driven variance reduction}
The diagnostics support the same mechanism as the entropic bottleneck rather than
a coverage-reshaping one. Final reaction-coordinate region fractions and the
fraction of effective ABF samples in the transition window are nearly identical
across ABF, FR estimated, FR uniform and FR oracle: FR does \emph{not} converge to
a different stationary $z$-distribution. Instead the gain appears immediately after
FR activates and is sustained for the rest of the run. \Cref{fig:wca_mechanism}
shows the interpretation: birth--death acts as a \emph{variance-reduction
resampler for the ABF estimator}. Replicas stuck in over-represented
configurations are replaced by whole-configuration clones of under-represented
ones, raising the number of effectively-independent force samples per $z$-bin even
when the marginal looks unchanged. Because the \emph{entire} WCA configuration is
copied (\cref{rem:manybody}), clones inject decorrelated solvent arrangements,
accelerating convergence of $\Fhat$ without requiring a different stationary
coverage --- consistent with $\text{uniform}\approx\text{estimated}\approx
\text{oracle}$.

\begin{figure}[t]
  \centering
  \includegraphics[width=0.92\textwidth]{fig_wca_04_mechanism.png}
  \caption{Mechanism. (a) Effective ABF sample count $N_{\mathrm{eff}}(z)$ is
    raised by FR in the transition window; (b) per-bin mean-force error is reduced
    there; (c) birth (clone) and death locations concentrate around the
    over/under-sampled regions; FR reshapes the \emph{sampling quality}, not the
    stationary marginal.}
  \label{fig:wca_mechanism}
\end{figure}

\subsection{Failure boundary}
Unlike the entropic bottleneck, the WCA dimer exhibits a clear failure boundary in
the FR rate (\cref{fig:wca_failure}). The rate ladder (a shorter $150$k-step,
$4$-seed stage) maps a clean safe regime: too gentle
(\texttt{fr\_rate}$\le0.05$) gives a small gain ($+19$--$29\%$); the sweet spot
(\texttt{fr\_rate}$\approx0.1$--$0.2$) gives the maximal $+58$--$64\%$ with healthy
ancestor ESS; the gain then \emph{erodes} sharply
(\texttt{fr\_rate}$=0.5$ falls to $+12\%$ with the error already climbing back
toward the ABF baseline) and finally \emph{reverses} for
\texttt{fr\_rate}$\ge1$, where the median final $\Ltwo(F)$ is worse than ABF
($-38\%$ at rate $1$, $-68\%$ at rate $2$) and no seed wins, as the ancestor ESS
collapses from $\sim\!337$ to $\sim\!20$. The reversal arrives even earlier under a
longer budget: the production aggressive configuration (\texttt{fr\_rate}$=0.5$ at
$250$k steps, $10$ seeds) is already net-harmful ($\WCAaggrGainPct\%$, winning only
$1/10$ seeds, ancestor ESS $\sim\!22$), because the extra steps give diversity loss
more time to dominate. Over-aggressive birth--death destroys the replica diversity
the mean-force estimator depends on, so the estimator degrades faster than coverage
improves. This is the predicted failure mode; the analytic $y$-channel of the
entropic bottleneck masks it (\cref{sec:case_eb}) because cloned walkers
re-equilibrate instantly, whereas a dense solvent does not.

\begin{figure}[t]
  \centering
  \includegraphics[width=0.92\textwidth]{fig_wca_05_failure.png}
  \caption{FR-rate ladder ($150$k steps, $4$ seeds). Final $\Ltwo(F)$ (left) is
    minimised in the $0.1$--$0.2$ sweet spot, climbs back toward the ABF baseline
    (dashed) by \texttt{fr\_rate}$=0.5$, and rises above it for
    \texttt{fr\_rate}$\ge1$; ancestor ESS (right) collapses as the rate grows,
    pinpointing diversity loss as the failure mechanism.}
  \label{fig:wca_failure}
\end{figure}

\subsection{Difficulty dependence}
Varying one difficulty axis at a time (\cref{fig:wca_difficulty}) shows the gain
is robust and, importantly, \emph{grows with difficulty}. With \emph{budget}, FR
wins at every $n_{\mathrm{steps}}$ and the relative advantage grows rather than
washes out ($+\WCAbudgetLoGain\%$ at $50$k to $+\WCAbudgetHiGain\%$ at $400$k):
ABF's error plateaus around $0.08$ after $\sim150$k while FR keeps descending, so
FR is not merely an early-time accelerator that ABF later catches. With
\emph{replica count}, the gain is roughly flat ($+51$--$59\%$ from $N=256$ to
$2048$) --- FR helps even at $N=256$. With \emph{solvent crowding}, the densest
bath ($a=1.35$) makes the entropic barrier severe and ABF $\sim4\times$ worse in
absolute error; FR still wins $5/5$ and delivers by far the largest absolute error
reduction ($\WCAdenseAbfErr\to\WCAdenseFrErr$). In the units a practitioner cares
about --- absolute free-energy error removed --- the crowded regime is where FR
helps most.

\begin{figure}[t]
  \centering
  \includegraphics[width=0.95\textwidth]{fig_wca_06_difficulty.png}
  \caption{Difficulty dependence (matched-seed gain over ABF). (a) Budget: the
    relative gain grows with $n_{\mathrm{steps}}$ as ABF plateaus and FR keeps
    improving. (b) Replica count: gain is roughly flat. (c) Solvent crowding:
    smaller $a$ (denser) yields the largest absolute error reduction. FR wins
    $5/5$ seeds in every cell.}
  \label{fig:wca_difficulty}
\end{figure}

\subsection{Case finding}
On the condensed-phase WCA dimer, ABF+FR with an online estimated target reliably
and substantially improves free-energy accuracy: $+\WCAtunedGainPct\%$ in median
final $\Ltwo(F)$ at the tuned operating point, winning $\WCAtunedWin/10$ seeds,
with the gain concentrated in the free energy (not the mean force). The mechanism
is balanced-resampling variance reduction after warm-up; the precise target shape
matters little ($\text{uniform}\approx\text{estimated}\approx\text{oracle}$, and
the estimated target already reaches the oracle's accuracy). The method has a clear
safe regime, \texttt{fr\_rate}$\approx0.1$--$0.2$, and a real failure boundary:
too-aggressive birth--death reverses the benefit as diversity collapses. The
advantage \emph{grows with simulation budget}, is insensitive to replica count, and
is largest in the densest, hardest solvent --- exactly the starved regime the
variance-reduction mechanism predicts.
